\documentclass[11pt]{article}

\usepackage[T1]{fontenc}
\usepackage[utf8]{inputenc}
\usepackage{amsmath,amsthm,mathtools}
\usepackage{newtxtext,newtxmath}
\usepackage{microtype}
\usepackage[letterpaper,margin=1in]{geometry}
\usepackage[dvipsnames]{xcolor}
\usepackage{booktabs,tabularx,array,longtable}
\usepackage{enumitem}
\usepackage{needspace}
\usepackage[most]{tcolorbox}
\usepackage{fancyhdr}
\usepackage{lastpage}
\usepackage{hyperref}
\usepackage{bookmark}

\definecolor{NPTnavy}{HTML}{18375F}
\definecolor{NPTteal}{HTML}{246B4F}
\definecolor{NPTamber}{HTML}{956000}
\definecolor{NPTred}{HTML}{9A3030}
\definecolor{NPTgray}{HTML}{66717E}
\definecolor{NPTbluewash}{HTML}{F2F6FA}
\definecolor{NPTgreenwash}{HTML}{F2F8F5}
\definecolor{NPTamberwash}{HTML}{FCF8EE}

\hypersetup{
  colorlinks=true,
  linkcolor=NPTnavy,
  citecolor=NPTteal,
  urlcolor=NPTteal,
  pdftitle={No Profitable Trade: From KKT to Revealed Preference},
  pdfauthor={Victor H. Aguiar},
  pdfsubject={Brainstorming paper and research agenda},
  pdfkeywords={KKT, revealed preference, Afriat, GARP, CMU, stochastic choice, continuum goods}
}

\setlength{\parindent}{0pt}
\setlength{\parskip}{5.5pt}
\setlist[itemize]{leftmargin=1.4em,itemsep=2.5pt,topsep=3pt}
\setlist[enumerate]{leftmargin=1.7em,itemsep=3pt,topsep=3pt}
\renewcommand{\arraystretch}{1.14}
\emergencystretch=1.5em

\pagestyle{fancy}
\fancyhf{}
\lhead{\small\color{NPTgray} Victor H. Aguiar}
\rhead{\small\color{NPTgray} No Profitable Trade}
\cfoot{\small\color{NPTgray} \thepage\ of \pageref*{LastPage}}
\renewcommand{\headrulewidth}{0.35pt}

\newcommand{\R}{\mathbb{R}}
\newcommand{\T}{\mathcal{T}}
\newcommand{\B}{\mathcal{B}}
\newcommand{\cG}{\mathcal{G}}
\newcommand{\cP}{\mathcal{P}}
\newcommand{\cR}{\mathcal{R}}
\newcommand{\ip}[2]{\left\langle #1,#2\right\rangle}
\newcommand{\pos}[1]{\left[#1\right]_{+}}
\newcommand{\argmax}{\operatorname*{arg\,max}}
\newcommand{\esssup}{\operatorname*{ess\,sup}}

\newtheorem{theorem}{Theorem}[section]
\newtheorem{proposition}[theorem]{Proposition}
\newtheorem{lemma}[theorem]{Lemma}
\newtheorem{corollary}[theorem]{Corollary}
\theoremstyle{definition}
\newtheorem{definition}[theorem]{Definition}
\newtheorem{example}[theorem]{Example}
\newtheorem{agenda}[theorem]{Research question}
\theoremstyle{remark}
\newtheorem{remark}[theorem]{Remark}

\newtcolorbox{statusbox}{
  enhanced,breakable,boxrule=0.7pt,arc=2pt,
  colframe=NPTamber,colback=NPTamberwash,
  left=9pt,right=9pt,top=7pt,bottom=7pt}
\newtcolorbox{keybox}{
  enhanced,breakable,boxrule=0.8pt,arc=2pt,
  colframe=NPTnavy,colback=NPTbluewash,
  left=10pt,right=10pt,top=8pt,bottom=8pt}
\newtcolorbox{knownbox}{
  enhanced,breakable,boxrule=0.7pt,arc=2pt,
  colframe=NPTteal,colback=NPTgreenwash,
  left=9pt,right=9pt,top=7pt,bottom=7pt}

\title{\vspace{-0.6in}\textbf{No Profitable Trade: From KKT to Revealed Preference}\\[0.35em]
  \large A Brainstorming Paper and Research Agenda}
\author{Victor H. Aguiar\\Simon Fraser University}
\date{October 9, 2026}

\begin{document}
\maketitle

\begin{statusbox}
\textbf{Status.} This document is a disciplined brainstorming paper, not a novelty claim and
not a finished submission.  It separates established mathematics from interpretive synthesis,
candidate diagnostics, and open research questions.  The organizing phrase ``no profitable
trade'' is deliberately economic; it is not asserted to be new terminology.  Any eventual paper
must add a systematic novelty search, sharpen the statistical objects, and choose one principal
theorem rather than pursue every extension at once.
\end{statusbox}

\begin{abstract}
This note develops a common language for constrained optimization and revealed preference.
For a proposed choice, exact regret asks whether another feasible bundle delivers higher utility;
its first-order counterpart is the support gap obtained by maximizing the current utility gradient
over the opportunity set.  For a concave objective, zero support gap is equivalent to global
optimality.  On a standard budget the gap has a closed form: it decomposes into unspent-income
loss and commodity-specific entry losses, reproducing all KKT conditions---including corners---
without requiring multipliers at the discovery stage.  Multipliers are recovered afterward as
dual certificates.

Across repeated budgets, the same geometry leads to Afriat inequalities.  Observation-specific
rescalings of prices serve as supporting utility gradients; GARP is the observable no-profitable-
cycle condition left after utility levels and marginal utilities of income are eliminated.  A
parallel blocking gap extends the idea to asymmetric preference functions and coherent
coalitional multi-utility representations, with an exact warning: no blocking is generally weaker
than preference maximization unless skew symmetry is imposed.  The note closes with proposed
stochastic-choice and continuum-good extensions, a formal-verification program, and an explicit
ledger distinguishing classical results from potentially new work.
\end{abstract}

\begin{keybox}
\centering
\textbf{One proposed conceptual spine}\par
\vspace{3pt}
\[
\boxed{
\begin{gathered}
\text{no improving feasible direction}
\;\Longrightarrow\;
\text{no profitable trade}\\[-1pt]
\Downarrow\\[-1pt]
\text{no profitable revealed-preference cycle}
\end{gathered}
}
\]
The first arrow is optimization within one opportunity set.  The second is an integrability
question across many opportunity sets.  They should be connected, but never conflated.
\end{keybox}

\tableofcontents

\section{The question and the levels of information}

The proposed program starts from one question:

\begin{quote}
Given a proposed or observed choice, can one exhibit a feasible trade, a revealed-preference
cycle, a blocking coalition, or a stochastic separating trial that demonstrates an improvement?
\end{quote}

The answer depends on what is observed.  It is useful to distinguish four levels.

\begin{enumerate}[label=\textbf{L\arabic*.}]
  \item \textbf{Known objective and one feasible set.}  We can compute exact regret and a
  first-order no-profitable-trade gap.
  \item \textbf{Known objective and many feasible sets.}  Every chosen point must pass the
  one-set test, and the gradients must arise from the same objective.
  \item \textbf{Unknown objective and finite demand observations.}  Supporting gradients and
  utility levels are latent.  Afriat's inequalities ask whether such objects exist; GARP is the
  corresponding observable cycle condition.
  \item \textbf{Population choice probabilities.}  A distribution of deterministic rational
  types is latent.  Convex separation replaces a single revealed-preference graph.
\end{enumerate}

This hierarchy prevents two recurring mistakes.  First, a numerical first-order gap is not
directly observable when utility is unknown.  Second, a revealed-preference cycle is not literally
an arbitrage or a sequence of trades that the consumer could execute through time.  It is a
logical inconsistency among the preference comparisons implied by choices, unless an additional
market mechanism gives the cycle a trading interpretation.

\section{One opportunity set: exact regret and the NPT gap}

Let $C\subset\R^L$ be nonempty, compact, and convex, and let $u$ be differentiable on a
neighborhood of $C$.  For $x\in C$, define the \emph{exact utility regret}
\begin{equation}
  \cR_u(x\mid C)
  :=\max_{y\in C}\{u(y)-u(x)\},
  \label{eq:exact-regret}
\end{equation}
and the \emph{linearized no-profitable-trade gap}
\begin{equation}
  \cG_u(x\mid C)
  :=\max_{y\in C}\ip{\nabla u(x)}{y-x}
  =\sigma_C(\nabla u(x))-\ip{\nabla u(x)}{x},
  \label{eq:npt-gap}
\end{equation}
where $\sigma_C(g):=\max_{y\in C}\ip{g}{y}$ is the support function of $C$.
The maximizing $y$ in \eqref{eq:npt-gap} is the best trade according to current marginal
utilities.  In optimization, this object is the variational-inequality or Frank--Wolfe gap
\cite{FrankWolfe1956,Jaggi2013}; the proposed contribution is to make its economic content
central and to connect it systematically to revealed preference.

\begin{proposition}[Exact regret is bounded by the NPT gap]
\label{prop:regret-gap}
Suppose $u$ is concave and differentiable.  For every $x\in C$,
\[
  0\le \cR_u(x\mid C)\le \cG_u(x\mid C).
\]
Moreover, the following statements are equivalent:
\begin{enumerate}[label=(\roman*)]
  \item $x\in\argmax_{y\in C}u(y)$;
  \item $\cR_u(x\mid C)=0$;
  \item $\cG_u(x\mid C)=0$;
  \item $\ip{\nabla u(x)}{y-x}\le0$ for every $y\in C$.
\end{enumerate}
\end{proposition}

\begin{proof}
Because $x\in C$, both maxima are at least zero.  Concavity gives
\[
  u(y)-u(x)\le\ip{\nabla u(x)}{y-x}
  \quad\text{for every }y\in C,
\]
and hence $\cR_u\le\cG_u$.  Thus $\cG_u=0$ implies global optimality.  Conversely,
if $x$ is globally optimal, then for each $y\in C$ the function
$\varphi(t)=u(x+t(y-x))$ is maximized at $t=0$ over $[0,1]$.  Therefore
$\varphi'(0)=\ip{\nabla u(x)}{y-x}\le0$.  Taking the maximum over $y$ proves
$\cG_u=0$.
\end{proof}

The proof makes the economics transparent.  Concavity converts the linearized value of the
best feasible trade into a valid upper bound on the value of the best actual deviation.  Without
concavity, $\cG_u=0$ remains a first-order stationarity condition but need not establish a global
maximum.

\begin{remark}[A positive gap certifies a direction, not necessarily a full jump]
If $\cG_u(x\mid C)>0$ and $y^\ell$ solves the support problem, then
$d=y^\ell-x$ has positive directional derivative.  Convexity of $C$ and differentiability of
$u$ imply that $x+\gamma d$ is feasible and improves utility for every sufficiently small
$\gamma>0$.  Curvature may nevertheless make the full jump from $x$ to $y^\ell$ unattractive.
A line search or a sufficiently small step converts the linear certificate into an actual
profitable trade.
\end{remark}

\subsection{Nonsmooth utilities}

For a concave function, let $\partial^+u(x)$ denote its superdifferential.  Given
$g\in\partial^+u(x)$, define
\[
  \cG(x;g,C):=\sigma_C(g)-\ip{g}{x}.
\]
The supergradient inequality immediately yields
\begin{equation}
  \cR_u(x\mid C)\le\cG(x;g,C).
  \label{eq:supergradient-certificate}
\end{equation}
Thus any $g\in\partial^+u(x)$ with $\cG(x;g,C)=0$ is a global optimality certificate.
This version is the right bridge to Afriat's piecewise-affine rationalizing utility, which is
typically nondifferentiable.

\subsection{Ordinal content and its limit}

Zero is robust, magnitude is not.  If $\widetilde u=\phi\circ u$ and $\phi$ is differentiable
with $\phi'(u(x))>0$, then
\[
  \nabla\widetilde u(x)=\phi'(u(x))\nabla u(x),
  \qquad
  \cG_{\widetilde u}(x\mid C)=\phi'(u(x))\cG_u(x\mid C).
\]
Consequently, the statement ``there is no profitable first-order trade'' is locally ordinally
invariant, but its utility-unit magnitude is not.  A numerical NPT gap must therefore be normalized
or tied to a specified cardinal representation before it is interpreted as a welfare quantity.

\section{A single budget: a closed form and KKT recovery}

Consider the standard budget
\[
  B(p,m):=\{y\in\R_+^L:p\cdot y\le m\},
  \qquad p\in\R_{++}^L,\quad m>0.
\]
Fix a feasible $x$ and write $g=\nabla u(x)$.  Define the largest positive marginal
utility per dollar
\begin{equation}
  \rho(x):=\max\left\{0,\max_{1\le i\le L}\frac{g_i}{p_i}\right\}.
  \label{eq:rho}
\end{equation}

\begin{theorem}[One-budget NPT formula]
\label{thm:budget-gap}
For every feasible $x$,
\begin{align}
  \cG_u(x\mid B(p,m))
  &=m\rho(x)-g\cdot x
  \label{eq:budget-gap}\displaybreak[0]\\
  &=\rho(x)\,[m-p\cdot x]
    +\sum_{i=1}^L p_ix_i
      \left[\rho(x)-\frac{g_i}{p_i}\right].
  \label{eq:budget-decomposition}
\end{align}
If $u$ is concave, $x$ is globally optimal if and only if this expression is zero.
Equivalently,
\begin{align}
  &m-p\cdot x\ge0,\qquad
    \rho\ge0,\qquad
    \rho(m-p\cdot x)=0,
  \label{eq:npt-budget-comp}\\
  &g_i-\rho p_i\le0,\qquad
    x_i(g_i-\rho p_i)=0
    \quad(i=1,\ldots,L).
  \label{eq:npt-entry-comp}
\end{align}
\end{theorem}

\begin{proof}
The linear support problem is
\[
  \max_{y\ge0,\,p\cdot y\le m}g\cdot y.
\]
If all $g_i\le0$, choosing $y=0$ gives value zero.  Otherwise, all expenditure is placed
on a good attaining the highest positive ratio $g_i/p_i$, so the value is $m\rho$.
Subtracting $g\cdot x$ yields \eqref{eq:budget-gap}; adding and subtracting
$\rho p\cdot x$ yields \eqref{eq:budget-decomposition}.  Every term in the latter is
nonnegative by feasibility and the definition of $\rho$.  Hence the sum is zero exactly when
all complementary products in \eqref{eq:npt-budget-comp}--\eqref{eq:npt-entry-comp} vanish.
Global optimality follows from Proposition~\ref{prop:regret-gap}.
\end{proof}

The formula turns the KKT conditions into an accounting identity.  Set
\begin{equation}
  \lambda:=\rho,
  \qquad
  \mu_i:=\rho p_i-g_i\ge0.
  \label{eq:recover-multipliers}
\end{equation}
Then $g-\lambda p+\mu=0$, while \eqref{eq:budget-decomposition} is precisely
\[
  \lambda(m-p\cdot x)+\mu\cdot x.
\]
For a generic candidate, $(\rho,\mu)$ is the optimal dual certificate for the \emph{linearized}
support problem.  Thus the NPT calculation can discover and screen an allocation without first
guessing multipliers.  When the gap is zero, concavity makes $x$ optimal for the original utility
problem and the same pair becomes an original-problem KKT certificate: $\rho$ is a shadow-income
supergradient and $\mu_i$ is the exclusion margin for a zero-consumption good.

\begin{corollary}[Interior equalization and corner screening]
If $\rho>0$, zero NPT gap requires the budget to bind.  Every consumed good satisfies
\[
  \frac{\partial_i u(x)}{p_i}=\rho,
\]
whereas every excluded good satisfies
\[
  \frac{\partial_i u(x)}{p_i}\le\rho.
\]
An excluded good with a strictly larger ratio is an immediate profitable-entry certificate.
\end{corollary}

\begin{example}[A corner detected before writing the Lagrangian]
\label{ex:corner}
Let
\[
  u(x_1,x_2)=a_1\log(1+x_1)+a_2x_2,
  \qquad a_1,a_2>0,
\]
and consider the candidate $\bar x=(0,m/p_2)$.  At this candidate,
\[
  g=(a_1,a_2),\qquad
  \rho=\max\left\{\frac{a_1}{p_1},\frac{a_2}{p_2}\right\},
\]
so
\[
  \cG_u(\bar x\mid B)
  =m\left[\frac{a_1}{p_1}-\frac{a_2}{p_2}\right]_+.
\]
The all-good-2 corner is globally optimal exactly when $a_1/p_1\le a_2/p_2$.  If the
inequality fails, the gap identifies good 1 as the profitable entrant.  The actual utility gain
from reoptimization is no larger than this linearized bound.
\end{example}

\begin{example}[Cobb--Douglas as an exact zero-gap audit]
For $u(x)=\alpha\log x_1+(1-\alpha)\log x_2$, the familiar demand
\[
  x_1^*=\frac{\alpha m}{p_1},
  \qquad
  x_2^*=\frac{(1-\alpha)m}{p_2}
\]
has
\[
  \frac{\partial_1u(x^*)}{p_1}
  =\frac{\partial_2u(x^*)}{p_2}
  =\frac1m.
\]
Hence $\rho=1/m$, $g\cdot x^*=1$, and $m\rho-g\cdot x^*=0$.  This one-line
calculation simultaneously checks budget exhaustion, equalized bang per buck, and global
optimality.  Here utility is understood on the positive orthant (or as an extended-value utility
equal to $-\infty$ on the axes); the gradient audit is applied at the strictly positive optimizer.
\end{example}

\subsection{A candidate scale-free diagnostic}

When $g\ge0$ and $\rho>0$, define
\begin{equation}
  \widehat\cG_u(x\mid B)
  :=\frac{\cG_u(x\mid B)}{m\rho}
  =1-\frac{g\cdot x}{m\rho}\in[0,1].
  \label{eq:normalized-gap}
\end{equation}
This quantity is invariant to multiplying $g$ by a positive scalar and therefore to a smooth
increasing transformation of utility with strictly positive derivative at $u(x)$.  It is a promising classroom diagnostic:
it reports the share of maximal current marginal value lost to slack income and misallocation.
It is \emph{not yet} a revealed-preference efficiency index or a cardinal welfare measure.
Its relationship to the critical cost efficiency index and other data-only measures is an open
question, not a theorem asserted here.

\section{Many linear constraints: multipliers as the dual of trade search}

Let
\[
  C=\{y\in\R_+^L:Ay\le b\}
\]
be nonempty and compact.  For a proposed $x\in C$ and a marginal-value vector $g$, linear
programming duality gives the exact identity
\begin{align}
  \cG(x;g,C)
  &=\max_{y\in C}g\cdot(y-x)
  \label{eq:poly-primal}\\
  &=\min_{\lambda\ge0:\,A^\top\lambda\ge g}
       \{b^\top\lambda-g^\top x\}
  \label{eq:poly-dual}\\
  &=\min_{\lambda\ge0:\,A^\top\lambda\ge g}
       \left\{\lambda^\top(b-Ax)
       +(A^\top\lambda-g)^\top x\right\}.
  \label{eq:poly-decomp}
\end{align}
Writing $\mu=A^\top\lambda-g\ge0$, a zero gap is equivalent to the existence of prices
of the constraints such that
\[
  g-A^\top\lambda+\mu=0,\qquad
  \lambda_j(b_j-A_jx)=0,\qquad
  \mu_ix_i=0.
\]
These are the KKT conditions under the maximization convention
$L=u(x)+\lambda^\top(b-Ax)+\mu^\top x$ \cite{KuhnTucker1951}.

This duality clarifies what a ``multiplier-free KKT'' can and cannot mean.  The primal support
problem avoids multipliers as a method of discovery: search directly for the best feasible trade.
The dual variables remain mathematically valuable because they compress the proof that no trade
is profitable.  Eliminating them from the student's first calculation does not eliminate duality.

\subsection{Sparse local tests}

Let $T_C(x)$ be the tangent cone of $C$ at $x$.  The first-order NPT condition is
\[
  g\cdot d\le0\qquad\text{for every }d\in T_C(x).
\]
For a polyhedron, it is enough to check a generating set of extreme rays of the tangent cone.
This suggests a sparse ``elementary trade'' implementation: pairwise swaps for a single budget,
network cycles for flow constraints, and circuit directions for a general matrix $A$.  The
underlying convex geometry is standard; the research task is to identify economically meaningful
generators and minimal certificates for important classes of models.

\section{From NPT to revealed preference}

Consider a finite demand dataset
\[
  \mathcal D=\{(p^t,x^t)\}_{t\in\T},
  \qquad p^t\in\R_{++}^L,\quad x^t\in\R_+^L\setminus\{0\},
\]
with observed wealth $m^t=p^t\cdot x^t$ and budget $B_t=B(p^t,m^t)$.
Bundle $x^t$ is directly revealed preferred to $x^s$ when
\[
  p^t\cdot x^s\le p^t\cdot x^t.
\]
GARP rules out a directed affordability cycle containing only weak edges and at least one strict
edge.  Equivalently, it rules out $x^t$ being indirectly revealed preferred to $x^s$ while
$x^s$ is strictly directly revealed preferred to $x^t$ \cite{Afriat1967,Varian1982}.

\subsection{The graph intuition}

Suppose first that one common, locally nonsatiated utility function rationalizes the data.  Every
outgoing affordability edge $t\to s$ implies
\[
  u(x^t)\ge u(x^s).
\]
Transitivity of numerical utility propagates this inequality along every directed path.  A strict
affordability edge, together with local nonsatiation, produces a strict preference comparison.
Therefore a strict directed cycle is impossible.

This yields the useful but carefully qualified slogan:
\begin{quote}
\textbf{NPT is no improving edge within a known opportunity set.  GARP is no inconsistent
cycle across observed opportunity sets when the objective is unknown.}
\end{quote}

\subsection{Afriat inequalities as compatible local certificates}

Afriat's theorem states that GARP is equivalent to the existence of numbers $U_t\in\R$ and
$\lambda_t>0$ satisfying
\begin{equation}
  U_s\le U_t+\lambda_t p^t\cdot(x^s-x^t)
  \qquad\text{for all }s,t\in\T.
  \label{eq:afriat}
\end{equation}
The piecewise-affine function
\begin{equation}
  u_A(x):=\min_{t\in\T}
  \left\{U_t+\lambda_t p^t\cdot(x-x^t)\right\}
  \label{eq:afriat-utility}
\end{equation}
is concave, continuous, strictly increasing, and rationalizes the data.  At $x^t$ the $t$th plane is
active, so
\begin{equation}
  q_t:=\lambda_t p^t\in\partial^+u_A(x^t).
  \label{eq:supporting-price}
\end{equation}
Moreover,
\[
  \cG(x^t;q_t,B_t)
  =\max_{y\in B_t}q_t\cdot(y-x^t)
  =0.
\]
Thus each observation receives a zero-NPT certificate, while inequalities
\eqref{eq:afriat} ensure that all certificates are supporting covectors of one concave potential.

\begin{theorem}[Afriat--NPT bridge: a restatement of known results]
\label{thm:afriat-npt}
For a finite Walrasian demand dataset, the following statements are equivalent:
\begin{enumerate}[label=(\roman*)]
  \item the dataset satisfies GARP;
  \item the Afriat inequalities \eqref{eq:afriat} have a solution with $\lambda_t>0$;
  \item the data are rationalizable by a continuous, strictly increasing, concave utility function;
  \item there exist a concave strictly increasing function $u$ and positive rescaled prices
  $q_t=\lambda_tp^t$ such that
  \[
    q_t\in\partial^+u(x^t)
    \quad\text{and}\quad
    \cG(x^t;q_t,B_t)=0
    \quad\text{for every }t;
  \]
  \item there are scalars $\lambda_t>0$ such that the graph
  $\{(x^t,q_t)\}_{t\in\T}$ is cyclically supermonotone:
  \begin{equation}
    \sum_{k=1}^K
    q_{t_k}\cdot(x^{t_{k+1}}-x^{t_k})\ge0,
    \qquad t_{K+1}=t_1,
    \label{eq:cyclic-supermonotone}
  \end{equation}
  for every finite cycle of observations.
\end{enumerate}
\end{theorem}

Statements (i)--(iii) are Afriat's theorem.  The equivalence with (iv) is the supporting-
hyperplane interpretation above.  Summing the supergradient inequalities around a cycle gives
(v).  Conversely, for fixed $q_t$ the Afriat system is the complete directed difference-constraint
system
\[
  U_s-U_t\le q_t\cdot(x^s-x^t).
\]
Such a system is feasible exactly when every directed cycle has nonnegative total weight.
The resulting $U_t$ and the minimum-of-affine construction \eqref{eq:afriat-utility} produce
the concave potential.  Finally, $q_t=\lambda_tp^t$ and
$p^t\cdot x^t=m^t$ imply that $q_t$ supports $B_t$ at $x^t$, hence its local support gap is
zero.  This is the finite-data form of Rockafellar's cyclic-monotonicity
construction \cite{Rockafellar1966}.  The word \emph{supermonotone} records the sign convention for concave
supergradients; $-q_t$ belongs to the usual convex cyclically monotone formulation.

\begin{keybox}
\textbf{Interpretation of $\lambda_t$.}  Price is not itself marginal utility.  The vector
$\lambda_t p^t$ is the supporting utility covector, and $\lambda_t$ is the observation-specific
marginal utility of income.  Raw price vectors generally need not be cyclically monotone because
income effects and cardinal utility scales vary across observations.  In quasilinear environments
with an interior numeraire, the normalization often fixes $\lambda_t=1$; this is why unscaled
cyclic monotonicity plays a direct role in quasilinear rationalizability \cite{Rochet1987}.
\end{keybox}

\subsection{Why corners do not break the bridge}

Suppose a differentiable concave $u$ rationalizes the observations and the KKT conditions at
observation $t$ are
\begin{equation}
  \nabla u(x^t)=\lambda_tp^t-\mu_t,
  \qquad \mu_t\ge0,
  \qquad \mu_t\cdot x^t=0.
  \label{eq:kkt-corner}
\end{equation}
For every observed $x^s$, concavity yields
\begin{align}
  u(x^s)-u(x^t)
  &\le \nabla u(x^t)\cdot(x^s-x^t)\\
  &=\lambda_tp^t\cdot(x^s-x^t)
    -\mu_t\cdot x^s+\mu_t\cdot x^t\\
  &\le\lambda_tp^t\cdot(x^s-x^t).
  \label{eq:corner-afriat}
\end{align}
This is the Afriat inequality.  The commodity-specific nonnegativity multipliers disappear from
the observable inequality because $-\mu_t\cdot x^s\le0$ and $\mu_t\cdot x^t=0$.
Corners therefore strengthen rather than obstruct the certificate.

\section{Preference functions and coalitional blocking}

Let $r:X\times X\to\R$ be a reflexive preference function normalized by $r(x,x)=0$,
where $r(x,y)\ge0$ means that $x$ is weakly preferred to $y$.  Standard utility is the special
case $r(x,y)=u(x)-u(y)$.  For a feasible set $B$ and $x\in B$, define two gaps:
\begin{align}
  \cP_r(x\mid B)
  &:=\sup_{y\in B}\pos{-r(x,y)},
  &&\text{preference-maximization gap},
  \label{eq:preference-gap}\\
  \B_r(x\mid B)
  &:=\sup_{y\in B}r(y,x),
  &&\text{blocking gap}.
  \label{eq:blocking-gap}
\end{align}
Reflexivity makes both gaps nonnegative.  The first is zero exactly when
$r(x,y)\ge0$ for every $y\in B$.  The second is zero exactly when no feasible $y$ satisfies
$r(y,x)>0$.

\begin{proposition}[Maximization, core, and skew symmetry]
\label{prop:blocking}
Suppose $r$ is asymmetric in the sense that
$r(x,y)\ge0$ implies $r(y,x)\le0$.
Then
\[
  \cP_r(x\mid B)=0\quad\Longrightarrow\quad \B_r(x\mid B)=0.
\]
If $r$ is skew-symmetric, $r(y,x)=-r(x,y)$, then
\[
  \cP_r(x\mid B)=\B_r(x\mid B),
\]
so preference maximization and no blocking are equivalent.
\end{proposition}

\begin{proof}
If $r(x,y)\ge0$ for every feasible $y$, asymmetry implies $r(y,x)\le0$.  Since
$r(x,x)=0$, the supremum in \eqref{eq:blocking-gap} is zero.  Under skew symmetry,
\[
  \sup_{y\in B}\pos{-r(x,y)}
  =\sup_{y\in B}\pos{r(y,x)}
  =\sup_{y\in B}r(y,x),
\]
where the last equality again uses the available zero at $y=x$.
\end{proof}

This proposition records an essential caveat: with a merely asymmetric preference function,
a core point need not be a preference maximizer.  Treating the blocking gap as exact regret is
valid under skew symmetry, but generally false without it.

\subsection{Coalitional multi-utility}

For a transparent finite formulation, let $\Omega$ be a nonempty finite collection of nonempty
finite sets $U$ of utility functions.  A CMU preference function has the form
\begin{equation}
  r_\Omega(y,x)
  =\max_{U\in\Omega}\min_{u\in U}\{u(y)-u(x)\}.
  \label{eq:cmu}
\end{equation}
Its blocking gap is
\begin{equation}
  \B_{r_\Omega}(x\mid B)
  =\sup_{y\in B}\max_{U\in\Omega}
    \min_{u\in U}\{u(y)-u(x)\}.
  \label{eq:cmu-blocking-gap}
\end{equation}
On a compact budget with continuous component utilities, zero means that no coalition $U$ and
feasible bundle $y$ make every utility in $U$ strictly favor $y$ over $x$.  The compact CMU
formulation replaces finiteness by the corresponding compactness and continuity hypotheses.

Aguiar, Hjertstrand, Serrano, and Evren characterize WGARP through asymmetric preference
functions and coherent CMU representations \cite{AguiarEtAl2026}.  Their Section~6 studies
exactly the associated core problem
\[
  \max_{y\in B}r_\Omega(y,\bar x).
\]
For such a coherent CMU, at $\bar x=x^*$ a preference maximizer solves this problem; the converse holds under skew
symmetry.  Thus the NPT language is not a replacement for that theory.  It is a proposed
unifying interpretation and a route to quantitative, computational, and formal extensions.

\Needspace{8\baselineskip}
For classical utility, \eqref{eq:cmu-blocking-gap} collapses to
\[
  \B_r(x\mid B)
  =\sup_{y\in B}\{u(y)-u(x)\}
  =\cR_u(x\mid B).
\]
The hierarchy can therefore be stated without ambiguity:
\[
  \cP_{r_\Omega}(x\mid B)=0
  \;\Longrightarrow_{\text{asymmetry}}\;
  \B_{r_\Omega}(x\mid B)=0,
  \qquad
  \cP_{r_\Omega}=0
  \;\Longleftrightarrow_{\text{skew symmetry}}\;
  \B_{r_\Omega}=0.
\]
For the singleton classical representation $r(x,y)=u(x)-u(y)$, skew symmetry holds and both
objects equal exact utility regret.

The magnitudes of $\cP_r$ and $\B_r$ depend on the cardinal representation $r$ and on
normalizations of the component utilities.  Their zero sets are the robust objects.  Any empirical
``distance from WGARP'' based on these gaps must specify an admissible normalization class.

\section{Stochastic choice: a separating-trial gap}

For a finite stochastic-choice design and a fixed population whose distribution of latent
preference types is stable across the observed choice sets, stack choice probabilities in a vector
$\pi\in\R^J$.  Let $a^1,\ldots,a^R$ denote the incidence vectors generated by all deterministic
rational choice types admitted by the design, and set
\[
  P:=\operatorname{conv}\{a^1,\ldots,a^R\}.
\]
Random-utility rationalizability is equivalent to $\pi\in P$ in this finite formulation
\cite{McFadden2005,KitamuraStoye2018}.  Given a norm $\|\cdot\|$ and its dual
$\|\cdot\|_*$, define
\begin{equation}
  \cG_{\mathrm{stoch}}(\pi)
  :=\max_{\|w\|_*\le1}
  \left\{w\cdot\pi-\sigma_P(w)\right\},
  \qquad
  \sigma_P(w)=\max_r w\cdot a^r.
  \label{eq:stochastic-gap}
\end{equation}
Convex duality gives
\begin{equation}
  \cG_{\mathrm{stoch}}(\pi)
  =\inf_{q\in P}\|\pi-q\|.
  \label{eq:stochastic-distance}
\end{equation}
Hence the gap is zero exactly under rationalizability.  When it is positive, an optimizing
$w$ is a unit-dual-norm normal vector to a separating hyperplane: a stochastic
revealed-preference trial on which the
observed probabilities outperform every deterministic rational type.  This is the probabilistic
counterpart of a profitable trade.

Equation \eqref{eq:stochastic-distance} is standard convex separation, not a claimed new
characterization.  The research opportunities are economic and statistical:
\begin{itemize}
  \item choose a norm and normalization with a transparent behavioral interpretation;
  \item obtain sparse separating trials that can be explained to a student or decision maker;
  \item derive sampling theory and critical values for the maximized gap;
  \item study robustness to measurement error, building on stochastic revealed-preference methods
  such as Aguiar and Kashaev \cite{AguiarKashaev2021};
  \item extend the construction from GARP-rational types to WGARP/CMU types.
\end{itemize}

\section{A continuum of differentiated goods}

The continuum-good extension is motivated by product environments in which the finite list of
commodities is itself an approximation.  Let $(\Omega,\mathcal F,\nu)$ be a nontrivial,
$\sigma$-finite measure space of product varieties, and let
$p:\Omega\to(0,\infty)$ be measurable.  The commodity space is the weighted cone
$L^1_+(p\,d\nu)$.  For a finite positive expenditure level $0<m<\infty$, define
\[
  B(p,m)=\left\{x\in L^1_+(p\,d\nu):
  \int_\Omega p(\omega)x(\omega)\,d\nu(\omega)\le m\right\}.
\]
Suppose the derivative of utility at $x$ is the continuous linear functional
\[
  D u(x)[d]=\int_\Omega g(\omega)d(\omega)\,d\nu(\omega),
\]
where $g$ is measurable and $g/p\in L^\infty(\nu)$.  Set
\begin{equation}
  \rho
  :=\esssup_{\omega\in\Omega}
  \max\left\{0,\frac{g(\omega)}{p(\omega)}\right\}.
  \label{eq:continuum-rho}
\end{equation}
Redefining the support gap with a supremum on this noncompact density space, the support
calculation gives
\begin{align}
  \cG(x;g,B)
  &:=\sup_{y\in B(p,m)}\int_\Omega g(\omega)[y(\omega)-x(\omega)]\,d\nu(\omega)
  \nonumber\\
  &=m\rho-\int_\Omega g(\omega)x(\omega)\,d\nu(\omega)
  \label{eq:continuum-gap}\\
  &=\rho\left[m-\int_\Omega p(\omega)x(\omega)\,d\nu(\omega)\right]
  +\int_\Omega p(\omega)x(\omega)
  \left[\rho-\frac{g(\omega)}{p(\omega)}\right]d\nu(\omega).
  \label{eq:continuum-decomp}
\end{align}
The proof is the finite-good argument after the change of variable
$z(\omega)=p(\omega)y(\omega)$.  The support value is $m$ times the essential supremum
of the positive part of $g/p$.  Sigma-finiteness permits expenditure to be concentrated on
finite-measure subsets on which this ratio approaches its essential supremum.  For $m>0$ and
$\rho>0$, a density maximizer exists exactly when
$\{\omega:g(\omega)/p(\omega)=\rho\}$ has positive measure; otherwise only a concentrating
sequence approaches the supremum.  If $\rho=0$, the zero density attains the support value.

The functional-analytic support identity is standard.  A serious continuum consumer theory must
still settle the economically important issues:

\begin{enumerate}[label=\textbf{C\arabic*.}]
  \item \textbf{Choice space.}  Densities, finite measures, and indivisible product choices imply
  different feasible spaces and dual objects.
  \item \textbf{Differentiability.}  The appropriate Fr\'echet, G\^ateaux, or directional derivative
  depends on the topology and on whether utility includes taste for variety or bundle interactions.
  \item \textbf{Attainment.}  Bliss-point or discrete-choice consumers may choose atoms rather than
  densities, changing both compactness and the support problem.
  \item \textbf{Neighborhood substitution.}  Cross-price effects should be aggregated over sets of
  nearby products, not reported as an unmanageably large matrix indexed by individual varieties.
  \item \textbf{Discretization.}  A theorem is needed showing when finite catalogs and their NPT
  certificates converge to the continuum model.
\end{enumerate}

\begin{agenda}[Neighborhood revealed preference]
Let products occupy a metric characteristic space $(\Omega,d)$ and let $N_h(\omega)$ be an
$h$-neighborhood.  Can one define revealed substitution kernels and neighborhood cross-price
elasticities that (i) are identified from discretized demand, (ii) converge as the catalog becomes
dense, and (iii) admit sparse NPT witnesses localized to economically interpretable product
neighborhoods?
\end{agenda}

\section{A formal-verification program}

The framework is well suited to proof certificates because failure and success both admit finite
witnesses in the finite-dimensional cases.

\begin{enumerate}
  \item \textbf{One budget.}  Verify the identity
  \[
    \cG=m\rho-g\cdot x
    =\rho(m-p\cdot x)+\sum_i x_i(\rho p_i-g_i)
  \]
  and conclude zero from nonnegative terms and complementarity.
  \item \textbf{General polyhedra.}  Verify primal feasibility, dual feasibility, stationarity,
  complementary products, and the supporting-hyperplane global bound.
  \item \textbf{GARP failure.}  A finite list of observations, affordability inequalities, and one
  strict inequality is a complete cycle certificate.
  \item \textbf{GARP success.}  Exact rational numbers $(U_t,\lambda_t)$ satisfying all Afriat
  inequalities certify rationalizability; the minimum-of-affine construction supplies the
  mathematical bridge.
  \item \textbf{Stochastic failure.}  A rational vector $w$ whose separating inequality is positive
  certifies nonmembership in the rational-type polytope.
\end{enumerate}

Lean should certify the smallest exact algebraic claim that matters; it should not be described as
having established novelty, prose quality, or an infinite-dimensional analytic theorem that has not
been formalized.  A practical repository would separate: (a) data certificates; (b) algebraic
verification; and (c) the analytic assumptions connecting those certificates to economic meaning.

\section{What is known, what is synthesis, and what remains open}

\begingroup
\small
\begin{longtable}{>{\raggedright\arraybackslash}p{0.18\textwidth}
  >{\raggedright\arraybackslash}p{0.32\textwidth}
  >{\raggedright\arraybackslash}p{0.40\textwidth}}
\toprule
\textbf{Status} & \textbf{Object} & \textbf{Assessment}\\
\midrule
\endfirsthead
\toprule
\textbf{Status} & \textbf{Object} & \textbf{Assessment}\\
\midrule
\endhead
\midrule
\multicolumn{3}{r}{\footnotesize Continued on the next page}\\
\endfoot
\bottomrule
\endlastfoot
Established
& Concavity, variational inequalities, and the support/Frank--Wolfe gap
& Zero gap characterizes a concave maximum; the gap bounds objective regret.\\
\addlinespace
Established
& KKT and linear-programming duality
& Multipliers are dual certificates for the absence of a profitable linearized trade.\\
\addlinespace
Established
& Afriat inequalities, GARP, and concave rationalization
& These equivalences are classical; the NPT vocabulary is a restatement and teaching bridge.\\
\addlinespace
Established
& Cyclic monotonicity and integrability
& Rockafellar supplies the mathematical potential theorem; positive rescaled prices, not generally
raw prices, are the relevant covectors.\\
\addlinespace
Established
& WGARP and coherent CMU rationalization
& Aguiar et al. provide the rationalization and core problem; no-blocking is weaker than maximization
without skew symmetry.\\
\addlinespace
Established
& Finite random-utility polytopes and separation
& Rationalizability is polytope/cone membership; a support inequality separates failures.\\
\addlinespace
Interpretive synthesis
& ``NPT $\to$ no profitable cycle $\to$ rationalizability''
& Potentially valuable as a unified exposition, but not itself a theorem of novelty.\\
\addlinespace
Candidate diagnostic
& Normalized local NPT gap \eqref{eq:normalized-gap}
& Locally invariant to utility rescaling; its empirical and welfare content remains to be proved.\\
\addlinespace
Research agenda
& Quantitative CMU and $k$-acyclic blocking gaps
& Requires a principled normalization and theorems relating gap size to observable violations.\\
\addlinespace
Research agenda
& Sparse stochastic separating trials
& Requires statistical theory, computational guarantees, and comparison with existing ARSP tests.\\
\addlinespace
Research agenda
& Continuum goods and neighborhood substitution
& Requires a choice-space model, existence and discretization results, and economically meaningful
cross-price objects.\\
\addlinespace
Research agenda
& Machine-checked certificate library
& Finite algebraic certificates are feasible; analytic bridges must be stated separately and honestly.\\
\end{longtable}
\endgroup

\section{A focused paper sequence}

Trying to establish every extension in one paper would obscure the central insight.  A credible
sequence works backward from sharply stated deliverables.

\subsection*{Paper 1: the finite deterministic bridge}

The first paper should center on one theorem or diagnostic beyond the classical equivalences.
A disciplined target is:
\begin{quote}
Construct a normalized, representation-aware NPT diagnostic for finite consumer data; characterize
its primal profitable-trade and dual price-support forms; and establish its exact relationship to
Afriat efficiency and sparse revealed-preference cycles.
\end{quote}
The literature review must determine whether this object is already implicit in approximate
revealed-preference, variational-inequality, or inverse-optimization work.

\subsection*{Paper 2: CMU and bounded cycles}

Define normalized blocking gaps for coherent CMU representations, determine what can be identified
without fixing a cardinal representation, and relate the minimal gap to WGARP and $k$-acyclicity.
The core/maximization distinction must remain explicit throughout.

\subsection*{Paper 3: stochastic certificates}

Develop sparse separating trials, finite-sample inference, and robust versions that allow
measurement error.  The contribution should be statistical interpretability or computational
sparsity, not the already-known convex separation characterization.

\subsection*{Paper 4: product neighborhoods in a continuum catalog}

Specify a characteristics-based continuum model, prove convergence from finite catalogs, and define
neighborhood substitution objects.  This paper should connect the billion-product motivation to a
mathematically disciplined limit experiment rather than merely replacing sums with integrals.

\section{Implications for the textbook}

Even before the research agenda is resolved, the established portion supplies a clean teaching
sequence:

\begin{enumerate}
  \item Utility maximization introduces the model and exact regret.
  \item The NPT gap shows why marginal conditions, corners, and global concavity belong together.
  \item KKT multipliers appear as recovered prices certifying that every feasible trade loses.
  \item Repeated observations turn local support into an integrability problem.
  \item Afriat and GARP solve the inverse problem: whether local certificates can be glued into one
  utility function.
  \item WGARP and CMU weaken the gluing requirement from all cycles to pairwise consistency and
  coalitional no blocking.
  \item Stochastic choice replaces a chosen bundle by a distribution and a profitable trade by a
  separating trial.
\end{enumerate}

This ordering makes revealed preference an implication and inverse problem of utility maximization,
not an apparently unrelated collection of axioms.  It also permits one notation throughout:
$x$ for a proposed choice, $B$ for its opportunity set, $g$ or $q$ for a supporting marginal-value
covector, and a gap for the best admissible challenge.

\section{Conclusion}

The core mathematics is simple enough to be memorable:
\[
\text{optimality}
\quad\Longleftrightarrow\quad
\max_{y\in C}\ip{\nabla u(x)}{y-x}=0
\]
for differentiable concave utility.  On a budget, the maximum has a closed form and decomposes into
the familiar KKT complementary products.  Across budgets, positive rescalings of price vectors are
supporting utility covectors; Afriat inequalities test whether they can be integrated into a single
concave potential, and GARP is the observable cycle restriction.  For nontransitive preference
functions, a blocking gap survives, but its equivalence to maximization requires skew symmetry.

The potential research contribution is not to rename KKT, Frank--Wolfe, Afriat, or Rockafellar.
It is to build an economics-native theory of certificates---local trades, global cycles, blocking
coalitions, stochastic trials, and continuum neighborhoods---with transparent normalization,
sparse witnesses, and formal verification.  The next step is to select one quantitative theorem
for the first paper and subject it to a dedicated novelty audit.

\begin{thebibliography}{99}

\bibitem{Afriat1967}
Afriat, S. N. (1967).
``The Construction of Utility Functions from Expenditure Data.''
\emph{International Economic Review} 8(1), 67--77.
\href{https://doi.org/10.2307/2525382}{doi:10.2307/2525382}.

\bibitem{AguiarEtAl2026}
Aguiar, V. H., P. Hjertstrand, R. Serrano, and O. Evren (2026).
``A Rationalization of the Weak Axiom of Revealed Preference.''
Working paper, arXiv:1906.00296, version 2 (manuscript dated November 2025;
arXiv revision posted August 12, 2026).
\url{https://arxiv.org/abs/1906.00296}.

\bibitem{AguiarKashaev2021}
Aguiar, V. H. and N. Kashaev (2021).
``Stochastic Revealed Preferences with Measurement Error.''
\emph{Review of Economic Studies} 88(4), 2042--2093.
\href{https://doi.org/10.1093/restud/rdaa067}{doi:10.1093/restud/rdaa067}.

\bibitem{FrankWolfe1956}
Frank, M. and P. Wolfe (1956).
``An Algorithm for Quadratic Programming.''
\emph{Naval Research Logistics Quarterly} 3(1--2), 95--110.
\href{https://doi.org/10.1002/nav.3800030109}{doi:10.1002/nav.3800030109}.

\bibitem{Jaggi2013}
Jaggi, M. (2013).
``Revisiting Frank--Wolfe: Projection-Free Sparse Convex Optimization.''
In \emph{Proceedings of the 30th International Conference on Machine Learning}, 427--435.

\bibitem{KitamuraStoye2018}
Kitamura, Y. and J. Stoye (2018).
``Nonparametric Analysis of Random Utility Models.''
\emph{Econometrica} 86(6), 1883--1909.
\href{https://doi.org/10.3982/ECTA14478}{doi:10.3982/ECTA14478}.

\bibitem{KuhnTucker1951}
Kuhn, H. W. and A. W. Tucker (1951).
``Nonlinear Programming.''
In J. Neyman (ed.), \emph{Proceedings of the Second Berkeley Symposium on Mathematical
Statistics and Probability}, 481--492. University of California Press.

\bibitem{McFadden2005}
McFadden, D. (2005).
``Revealed Stochastic Preference: A Synthesis.''
\emph{Economic Theory} 26(2), 245--264.
\href{https://doi.org/10.1007/s00199-004-0495-3}{doi:10.1007/s00199-004-0495-3}.

\bibitem{NishimuraOk2016}
Nishimura, H. and E. A. Ok (2016).
``Utility Representation of an Incomplete and Nontransitive Preference Relation.''
\emph{Journal of Economic Theory} 166, 164--185.
\href{https://doi.org/10.1016/j.jet.2016.07.002}{doi:10.1016/j.jet.2016.07.002}.

\bibitem{Rochet1987}
Rochet, J.-C. (1987).
``A Necessary and Sufficient Condition for Rationalizability in a Quasi-Linear Context.''
\emph{Journal of Mathematical Economics} 16(2), 191--200.
\href{https://doi.org/10.1016/0304-4068(87)90007-3}{doi:10.1016/0304-4068(87)90007-3}.

\bibitem{Rockafellar1966}
Rockafellar, R. T. (1966).
``Characterization of the Subdifferentials of Convex Functions.''
\emph{Pacific Journal of Mathematics} 17(3), 497--510.
\url{https://msp.org/pjm/1966/17-3/pjm-v17-n3-p08-p.pdf}.

\bibitem{Samuelson1938}
Samuelson, P. A. (1938).
``A Note on the Pure Theory of Consumer's Behaviour.''
\emph{Economica} 5(17), 61--71.

\bibitem{Varian1982}
Varian, H. R. (1982).
``The Nonparametric Approach to Demand Analysis.''
\emph{Econometrica} 50(4), 945--973.
\href{https://doi.org/10.2307/1912771}{doi:10.2307/1912771}.

\end{thebibliography}

\end{document}
